\begin{abstract}
Personalization is often assumed to increase trust in AI systems and encourage users to delegate decisions to algorithms. However, this preregistered experiment provides evidence that personalization can actually deter delegation. In a preregistered between-subjects experiment (N = 233), participants repeatedly chose whether to decide themselves or to delegate to an algorithm. In one treatment, the delegation option is an optimization-based algorithm selecting the option with the highest expected value. In the other treatment, it is a preference-based (personalized) algorithm selecting according to participants’ previously elicited preferences. We find no evidence that preference-based personalization increases delegation relative to expected-value optimization throughout the task even though the personalized algorithm should theoretically lead to higher expected utility for them. Instead, personalization deters delegation in two ways. First, at first contact participants are significantly less likely to delegate to the preference-based algorithm than to the optimization-based algorithm (66.9\% vs. 55.8\%, p = 0.0413); this difference is confined to the very first decision and disappeared immediately thereafter. 
Second, categorical non-delegation (never delegating) was substantially more prevalent under preference-based personalization (15.7\%) than under objective optimization (3.4\%). 
 Taken together, our findings challenge common assumptions about the benefits of personalization in AI systems and have important implications for the design of algorithmic agents. Designers should carefully consider when to personalize an algorithm and when to use objective algorithms, particularly in contexts where objective criteria are available and valued.

\textbf{Keywords}: AI delegation, personalization, algorithm appreciation, trust in AI, financial decision-making, preregistered experiment
\end{abstract}